\section{Introduction}

Do video world models learn physics, or do they memorize what physics usually looks like? Two
kinds of evidence are usually offered, and neither settles the question.

\emph{Behavioural} evidence asks whether generated futures obey physical laws. Controlled studies
find near-perfect in-distribution behaviour and failure out of distribution (OOD)
\citep{kang2025far}. Large video generators fail physics benchmarks while producing realistic
pixels \citep{motamed2026physicsiq,zhang2025morpheus,thozhiyoor2025slower}.

\emph{Representational} evidence asks whether physical variables can be decoded from internal
states \citep{joseph2026interpreting,tan2026latent,interno2026observer}. Decodable conserved
quantities can be causally inert \citep{liao2026objective}.

Behaviour tells us \emph{that} an OOD failure happens, not \emph{where}. Probes tell us that
information is present, not that it is used, and certainly not what an ideal learner would have
done.

This paper asks one sharp version of the question. A world model must infer gravity, a hidden
physical parameter, from a few context frames and then predict the future. When the true gravity
lies outside the training range, does the model fail because its belief state never represents an
out-of-range gravity (\emph{inference}), because its transition cannot carry one out
(\emph{execution}), or because its perception misreads the frames (\emph{perception})? And is the
failure a failure at all?

Our first contribution answers the last question in theory (Section~\ref{sec:theory}). Under
squared loss, the Bayes-optimal predictor's implied gravity is the posterior mean, which lies inside
the training range for every context. With noisy observations it converges to the boundary. Without
noise, the objective places no constraint at all on out-of-range contexts. Snapping back is
therefore what the loss asks for, and any extrapolation of the law has to come from an inductive
bias. We also show that low loss alone forces physical quantities to be decodable. A positive probe
is consequently weak evidence of anything beyond competence.

Our second contribution is a causal dissection of the failure. We build a controlled world of
balls and blocks with per-episode gravity, restitution and friction, in which materials are visible
and gravity is not. Three families of chunk world models share one dynamics architecture and differ
only in the training signal: pixel reconstruction (REC), latent prediction (JEPA) and ground-truth
positions (STATE). For each model we measure gravity three ways: the gravity its predictions
exhibit, the gravity a probe reads from its belief, and the gravity its predictions exhibit when
we \emph{write} a gravity into the belief. The write uses a direction estimated from exact
matched-state counterfactuals, made possible by time-reversing the simulator.

Because a linear write can fail for uninteresting reasons, we add \emph{nested inversion}:
\begin{itemize}
\item optimize the belief itself, which tests execution;
\item optimize the context latents fed to the frozen belief network, which tests reachability.
\end{itemize}
In both cases we fit only the first predicted frames and measure physics on held-out frames. An
anti-physics control (upward gravity) calibrates how flexible the decoder is.

Everything was pre-registered, including a first version of the study whose models never became
competent. We report it in Appendix~\ref{app:v1}, because the way it failed is itself informative.
